\section{Bitácora de reproducibilidad de las simulaciones preliminares}

Los resultados de la Tabla 1 y la Figura 1 se generan con un prototipo de producto
constante, una referencia CEX exógena, arbitraje acotado y una regla opcional de retiro
adaptativo de liquidez. La versión pública del código, el resumen de métricas, la figura
y el informe metodológico se encuentran en las rutas indicadas en la sección de
arquitectura. Para esta comprobación se ejecutaron 100 semillas por variante. E0--E2
son comparaciones mecánicas sin ruido intencional; E3 incluye innovaciones, tendencia y
ventas forzadas. Por ello, la repetición numérica de E0--E2 no debe confundirse con
independencia estadística.

El punto de partida reproducible es:

\begin{quote}
\texttt{python3 experiments/cex\_dex\_preliminary/cex\_dex\_preliminary.py}\\
\texttt{--output output/preliminary\_simulations --seeds 100}
\end{quote}

Los artefactos mínimos a conservar son: el archivo de configuración o sus parámetros
equivalentes, las semillas, la versión del intérprete, el resumen
\texttt{preliminary\_summary.csv}, las trayectorias cuando proceda y la figura
reproducida. Una réplica futura debe tratar las simulaciones actuales como selección de
mecanismos; las configuraciones de confirmación y cualquier conjunto empírico deben
separarse antes de informar métricas conclusivas.

\subsection*{Disponibilidad de código y datos}

El repositorio público de investigación se encuentra en
\href{https://github.com/edwardcalderon/Master-s-thesis}{Master-s-thesis}. Esta versión incluye el código,
la figura, el resumen de resultados y el informe de las simulaciones preliminares en las
rutas citadas en la sección de arquitectura. Los resultados actuales son sintéticos y
no se presentan como validación empírica.
